% ============================================================================
% SEC02.TEX - Section 5.2: Membuat Script Pergerakan Pesawat
% ============================================================================

\section{Membuat Script Pergerakan Pesawat}

Mari buat script agar pesawat bisa dikendalikan.

\begin{enumerate}
    \item Di folder \texttt{Scripts}, buat C\# Script baru bernama \texttt{PlayerController}.
    \item Attach script tersebut ke GameObject \texttt{Player}.
    \item Buka script dan tulis kode berikut:
\end{enumerate}

\begin{lstlisting}[language={[Sharp]C}, caption=PlayerController.cs (Movement)]
using UnityEngine;

public class PlayerController : MonoBehaviour
{
    public float moveSpeed = 5f;
    private Rigidbody2D rb;
    private Vector2 movement;

    void Start()
    {
        rb = GetComponent<Rigidbody2D>();
    }

    void Update()
    {
        // 1. Baca Input
        float moveX = Input.GetAxisRaw("Horizontal");
        float moveY = Input.GetAxisRaw("Vertical");
        
        movement = new Vector2(moveX, moveY).normalized;
    }

    void FixedUpdate()
    {
        // 2. Gerakkan Rigidbody
        rb.velocity = movement * moveSpeed;
    }
}
\end{lstlisting}

\subsection{Penjelasan Kode}
\begin{itemize}
    \item \texttt{GetComponent<Rigidbody2D>()}: Mengambil referensi komponen fisik pesawat.
    \item \texttt{Input.GetAxisRaw}: Membaca input keyboard secara instan (tanpa smoothing) agar respon lebih cepat.
    \item \texttt{FixedUpdate}: Tempat terbaik untuk memanipulasi fisika (Rigidbody).
\end{itemize}
